\chapter*{Introducción}
\addcontentsline{toc}{chapter}{Introducción}

El transporte público es la columna vertebral de la movilidad en las grandes urbes modernas. En la Zona Metropolitana de Guadalajara (ZMG), el crecimiento demográfico y la expansión de la mancha urbana han generado una red de transporte extensa y compleja, cuya eficiencia no depende únicamente de la cantidad de unidades o rutas, sino de la estructura subyacente que las conecta. Entender cómo se articulan estas conexiones es fundamental para diagnosticar problemas de movilidad, identificar vulnerabilidades y planificar mejoras estratégicas. En el contexto de la ZMG, el análisis de transporte se ha enfocado en métricas operativas aisladas (frecuencias, tiempos de recorrido) o en modelos macroscópicos de flujo \cite{inegi2024, iieg2024}. Sin embargo, estos enfoques a menudo pasan por alto la arquitectura topológica del sistema: ¿cómo está "tejida" la red?, ¿existen puntos críticos cuya falla colapsaría el servicio? ¿Se agrupan las rutas en comunidades funcionales desconectadas de la geografía administrativa?

% TODO [PDF p.9][Fernando Becerra][cat: referencias][prio: alta][resaltado: "enfoques previos"]: agregar o revisar referencias para sustentar esa comparación con trabajos previos.
Este trabajo aborda esas interrogantes desde la perspectiva de la \textbf{Teoría de Grafos}. A diferencia de enfoques previos que simplifican el sistema a un único modelo, esta investigación adopta una metodología dual basada en \textbf{espacios de representación}: el espacio de infraestructura (Espacio $L$), que modela la red física de paradas y tramos; y el espacio de servicio (Espacio $P$), que modela la red de transferencias y viajes directos desde la perspectiva del usuario.

Esta distinción permite un análisis más rico y matizado. Mientras que el Espacio $L$ nos habla de la cobertura territorial y la geometría de la red, el Espacio $P$ revela la eficiencia funcional, la necesidad de transbordos y la integración operativa. Utilizando los datos abiertos de la especificación GTFS (\textit{General Transit Feed Specification}) de la ZMG, se construyen y analizan ambos modelos matemáticos para caracterizar el estado actual del sistema.

\paragraph{Tesis Central}

% TODO [PDF p.9][Fernando Becerra][cat: metodologia][prio: alta][resaltado: "relativamente robusto"]: definir operativamente qué significa "robusto" y qué es "relativo" en este contexto.
Esta investigación sostiene y pone a prueba una proposición principal: \textbf{la ZMG combina una base física sensible a fallas con un servicio relativamente robusto, efecto explicado por la redundancia entre rutas.}

% TODO [PDF p.9][ernes][cat: metodologia][prio: alta][resaltado: "Planteamiento del Problema"]: formular una o dos preguntas explícitas de investigación, por ejemplo sobre robustez ante fallos dirigidos y nodos críticos cuya eliminación fragmenta la red.
\paragraph{Planteamiento del Problema}

La red de transporte de la ZMG ha evolucionado de manera orgánica. Si bien existen esfuerzos de planificación, carecemos de evaluaciones cuantitativas públicas que analicen la robustez estructural de la red actual frente a fallos. Estudios previos, como el de Barraza y Estrada (2021), han sentado un precedente importante al aplicar teoría de grafos a redes teóricas u ortogonales en la región. Sin embargo, existe la necesidad de aplicar este marco analítico a la red operativa real y actual, desglosando sus propiedades en función de la infraestructura física frente a la oferta de servicio, y explorando métricas avanzadas de modularidad y resiliencia que no han sido abordadas en profundidad localmente.

\paragraph{Objetivos}

El objetivo general de este trabajo es \textbf{modelar y analizar estructuralmente el sistema de transporte público de la Zona Metropolitana de Guadalajara mediante la Teoría de Grafos}, diferenciando entre la topología de la infraestructura y la topología del servicio para evaluar su eficiencia y robustez.

De este objetivo se desprenden los siguientes objetivos específicos:
\begin{itemize}
    \item Construir los grafos correspondientes al Espacio $L$ (infraestructura) y Espacio $P$ (servicio) a partir de los datos GTFS oficiales.
    \item Calcular y comparar métricas topológicas fundamentales (grado, intermediación, cercanía, coeficiente de agrupamiento) en ambos espacios para identificar nodos críticos.
    \item Analizar la estructura de comunidades de la red mediante algoritmos de modularidad, contrastando la organización geográfica con la funcional.
    \item Evaluar la robustez de la red simulando escenarios de fallos aleatorios y ataques dirigidos, cuantificando la degradación del servicio.
\end{itemize}

\paragraph{Estructura}

El documento se organiza en cinco capítulos que guían al lector desde los conceptos fundamentales hasta los hallazgos empíricos:

\begin{description}
    \item[Capítulo 1: Definiciones.] Establece el marco matemático necesario, introduciendo los conceptos de teoría de grafos, métricas de centralidad y propiedades de redes grafos que se utilizarán a lo largo del trabajo.
    
    \item[Capítulo 2: Análisis del transporte público.] Revisa el estado del arte en la aplicación de redes complejas al transporte. Se definen formalmente los espacios de representación ($L$, $P$, $C$, $B$) y se discute la literatura relevante, incluyendo los antecedentes locales.
    
    \item[Capítulo 3: Datos y Construcción.] Detalla la metodología de procesamiento de datos. Explica cómo se transforman los archivos GTFS crudos en grafos matemáticamente tratables, describiendo los procesos de limpieza y las decisiones de modelado tomadas.
    
    \item[Capítulo 4: Resultados.] Presenta el análisis cuantitativo de la red de la ZMG. Se contrasta explícitamente la fragilidad del Espacio $L$ frente a la resiliencia del Espacio $P$ mediante métricas, nodos centrales, comunidades y pruebas de robustez.
     
    \item[Capítulo 5: Conclusiones.] Sintetiza los hallazgos principales, verifica la tesis central de infraestructura frágil versus servicio resiliente, responde a las preguntas de investigación y propone líneas futuras de trabajo.
\end{description}

Esta investigación aspira a demostrar que las herramientas matemáticas de la teoría de grafos son instrumentos potentes para comprender y mejorar la realidad compleja de nuestros sistemas urbanos.
